\section{Exact quadratic support on polytopes}\label{app:polytope}

\subsection{Proof of Theorem~\ref{thm:polytope}}

Every candidate lies in $P$, so an empty $P$ has none. Assume $P\ne\emptyset$.
Among the global minimizers, which exist by compactness, choose $x^\ast$ such
that the face $G$ of $P$ with $x^\ast\in\relint G$ has the smallest dimension,
and let $V$ be the linear space parallel to $\aff G$. Because $x^\ast$
minimizes $q$ over $G$ and lies in its relative interior,
$v\T(Hx^\ast+c)=0$ and $v\T Hv\ge0$ for all $v\in V$.

Suppose $v\in V\setminus\{0\}$ satisfies $v\T Hv=0$. Then $q(x^\ast+tv)=q(x^\ast)$
for all $t$. The line $x^\ast+tv$ meets the bounded face $G$ in a segment that
contains $x^\ast$ in its relative interior, and an endpoint of the segment is
a global minimizer on a proper face of $G$, which contradicts the choice of
$x^\ast$. Hence $H$ is positive definite on $V$, or $V=\{0\}$.

Let $J$ be the set of rows active at $x^\ast$, so that
$\aff G=\{x:A_Jx=b_J\}$ and $V=\ker A_J$, and let $I\subseteq J$ index a basis
of the row space of $A_J$. Then $\abs I\le d$ and $\ker A_I=V$. Since
$Hx^\ast+c$ is orthogonal to $V$, it lies in the row space of $A_I$, so
$(x^\ast,\mu)$ solves~\eqref{eq:kkt-system} for some $\mu$. If $M_I(v,\nu)=0$,
then $A_Iv=0$, so $v\in V$, and $v\T Hv=-v\T A_I\T\nu=0$; hence $v=0$, and
$A_I\T\nu=0$ gives $\nu=0$ by independence. So $M_I$ is nonsingular, its unique
solution contains $x^\ast$, and $x^\ast$ is a candidate. All data are rational,
so the candidates and their values are rational. A set of dependent rows
never gives a nonsingular $M_I$, so enumerating independent sets suffices.
\qed

\subsection{Encoding lengths}

Multiply each row of $Ax\le b$ by the product of the denominators of its
entries, and multiply $q$ by the product $\sigma$ of the denominators of its
coefficients. Positive row scalings and $\sigma>0$ change neither $P$, nor the
nonsingularity of $M_I$, nor the $x$-part of its solution. Afterwards all row
entries are integers of absolute value at most $2^{\langle A\rangle}$ and all
entries of $\sigma H$ and $\sigma c$ are integers of absolute value at most
$2^{\langle q\rangle+1}$; write $\alpha=\langle A\rangle$ and
$\eta=\langle q\rangle+1$.

Let $\abs I=s\le d$ and $n=d+s$. Scale the last $s$ rows and the last $s$
columns of the integer matrix $[M_I\mid(-\sigma c,b_I)]$ by $2^{-\alpha}$. In
the scaled matrix, the first $d$ rows have entries of absolute value at most
$2^\eta$ and the others at most $1$. Hadamard's inequality applied to any
square submatrix of order at most $n$, and the undoing of the scaling, give
\[
\abs{\det N}\le n^{n/2}2^{d\eta+2s\alpha}\le B:=(2d)^d2^{d(\eta+2\alpha)}
\]
for every square submatrix $N$ of $[M_I\mid(-\sigma c,b_I)]$. By Cramer's rule a
candidate has coordinates $x_j=N_j/\Lambda$ with $\abs{N_j},\abs\Lambda\le B$.
Writing $\sigma q=\tilde c_0+\sum_j\tilde c_jx_j+\sum_{i\le j}\tilde Q_{ij}x_ix_j$
with integer coefficients bounded by $2^\eta$,
\[
\sigma q(x)=\frac{\Lambda^2\tilde c_0+\Lambda\sum_j\tilde c_jN_j+\sum_{i\le j}\tilde Q_{ij}N_iN_j}{\Lambda^2},
\]
a numerator with at most $(d+1)^2$ terms, each bounded by $2^\eta B^2$. So
$\min_Pq$ is a fraction whose numerator and denominator are bounded by
$(d+1)^22^{\eta}B^2$ and $2^{\eta}B^2$. Hence every candidate and the minimum
have bit length $O(d(\langle q\rangle+\langle A\rangle+\log d))$, for every $d$.

Bareiss's fraction-free elimination keeps every intermediate number equal to
a minor of the matrix above, and the products formed before each exact
division are bounded by $2B^2$. Feasibility tests and value comparisons are
integer comparisons of numbers bounded by $(d+1)2^\alpha B$ and
$(d+1)^22^\eta B^4$. The total work is $O(N(m,d)(d^3+md))$ arithmetic
operations on integers of bit length $O(d(\langle q\rangle+\langle A\rangle+\log d))$.
For fixed $d$, $N(m,d)\le(m+1)^d$, the number of arithmetic operations does
not depend on the sizes of the numbers, and the algorithm is strongly
polynomial.

\subsection{Proof of Proposition~\ref{prop:maxcut}}

For a graph $G=(V,E)$ let $q(x)=-\sum_{\{i,j\}\in E}(x_i+x_j-2x_ix_j)$. Since
$q$ is affine in each coordinate, moving coordinates one at a time to a
minimizing endpoint does not increase $q$, so some minimizer is binary. At a
binary point the summand of an edge equals one exactly when the edge crosses
the cut defined by $x$, so $\min q$ is minus the maximum cut size, and
unweighted MAX CUT is NP-hard \citep{Karp1972,GareyJohnsonStockmeyer1976}. The
coefficients are integers of absolute value at most $n$. For the second
statement, membership in coNP follows because a violating point can be taken
to be a candidate of Theorem~\ref{thm:polytope}, which has polynomial size.
For the hardness, $\min q$ is an integer for the quadratic $q$ of the
reduction, because $q$ has integer coefficients and a binary minimizer; hence
$\min q\le-W$ holds if and only if $q\ge-W+\tfrac12$ fails on $[0,1]^n$. \qed
